\subsection{When the margins certify a decision}
\label{app:certificate-validation}

We test whether valid envelopes can resolve candidate decisions, rather than
only contain observed errors. This is an analytic \textsc{known}-law check:
the reference laws and an initial-state box are supplied. It does not
calibrate a learned model or infer the box from partial observations.

Three DGUs form a chain, with unit capacitances, inductances, and resistances
in their physical units. Initially, the bus voltages are $(22,24,26)$ V,
the DGU currents are $(2,2,2)$ A, and the two line currents are zero; each
node draws a known constant 2 A load. The nine candidates have source
voltages
\[
 u_{c,s}=26\mathbf 1+c\mathbf 1+s(-1,0,1)^\top\ \mathrm V,
 \qquad c\in\{-2,0,2\},\quad s\in\{-1.5,0,1.5\}.
\]
Each candidate keeps the graph fixed and rolls out 100 steps of 0.01 s.
The maintained C-PHAST world-model interface calls its
\texttt{CompositionalPHAST} core with four substeps. DGU resistance in the
predictor is either the correct $1\,\Omega$ or a prescribed
$1.1\,\Omega$; no parameters are fitted. The coordinatewise initial radius
is $r\in\{0,0.001,0.01,0.05,0.2,1\}$ V for voltage coordinates and A for
current coordinates, shared by all candidates.

The score is the existing voltage-regulation/current-sharing objective,
averaged over samples $1{:}{N_{\mathrm{pred}}}$:
\[
 J(a)=\frac1{N_{\mathrm{pred}}}\sum_{k=1}^{{N_{\mathrm{pred}}}}\left[
 \left(\frac{\bar V_k-24\,\mathrm V}{1\,\mathrm V}\right)^2+
 \frac13\sum_{i=1}^3
 \left(\frac{I_{i,k}-\bar I_k}{1\,\mathrm A}\right)^2\right],
 \qquad \bar V_k=\frac13\sum_i V_{i,k},\quad
 \bar I_k=\frac13\sum_i I_{i,k}.
\]
The constraint limits the largest absolute sampled line current. We retain
all four prescribed diagnostic limits, $0.25$, $0.5$, $1$, and $2$ A;
these are not hardware ratings.

\paragraph{Enclosure.}
The independent reference dynamics are linear, $\dot x=Ax+Gu$, with
$x=(V_1,I_1,V_2,I_2,V_3,I_3,I_{e,1},I_{e,2})^\top$ and
$u=(-I_{\mathrm{load},1},V_{\mathrm{src},1},\ldots,-I_{\mathrm{load},3},V_{\mathrm{src},3})^\top$.
Here each load is 2 A and the three source-voltage entries are $u_{c,s}$;
the six-input forcing includes both, with units carried by $G$.
We enclose the resulting $14\times14$ augmented exponential
$\exp(\Delta t[\begin{smallmatrix}A&G\\0&0\end{smallmatrix}])$
using 50-digit outward-rounded interval arithmetic, a 40-term Taylor
polynomial, and the remainder
$e^{\|Z\|_\infty}\|Z\|_\infty^{41}/41!$, where
$Z=\Delta t[\begin{smallmatrix}A&G\\0&0\end{smallmatrix}]$.
Interval powers retain the signed initial-state sensitivity over the horizon.
If $[E_k]$ encloses the state-transition block and $[c_k(a)]$ the reference
trajectory from the box center, then every allowed initial state satisfies
\[
 x_{i,k}(a)\in[c_{i,k}(a)]+[-r,r]\sum_j|[E_{k,ij}]|.
\]
Interval evaluation of the score and peak current gives uniform consequence
enclosures. Their maximum deviations from the native predictor's scores give
$\Delta_J$ and $\Delta_g$, including the prescribed model mismatch and
discretization error. This direct enclosure instantiates the decision-margin
argument without requiring a contraction estimate for this electrical system.

\begin{table}[t]
\centering
\small
\caption{Analytic certificate informativeness at the prescribed 1 A sampled
line-current limit. Feasible, infeasible, and unresolved counts sum to nine;
rankings count strictly separated score envelopes among 36 unordered pairs.
Color distinguishes resolved feasible cases (green) from unresolved cases
(orange), not a performance ranking.}
\label{tab:certificate-informativeness}
\begin{tabular}{ccrrrr}
\toprule
Predictor $R$ ($\Omega$) & Radius $r$ & Feasible & Infeasible & Unresolved & Rankings \\
\midrule
1.0 & 0     & \goodcell{6} & 3 & 0 & 33/36 \\
1.0 & 0.001 & \goodcell{6} & 3 & 0 & 33/36 \\
1.0 & 0.01  & \goodcell{6} & 3 & 0 & 33/36 \\
1.0 & 0.05  & \goodcell{6} & 0 & \warncell{3} & 29/36 \\
1.0 & 0.2   & 0 & 0 & \warncell{9} & 0/36 \\
1.0 & 1     & 0 & 0 & \warncell{9} & 0/36 \\
\midrule
1.1 & 0     & \goodcell{6} & 3 & 0 & 32/36 \\
1.1 & 0.001 & \goodcell{6} & 3 & 0 & 32/36 \\
1.1 & 0.01  & \goodcell{6} & 3 & 0 & 30/36 \\
1.1 & 0.05  & \goodcell{6} & 0 & \warncell{3} & 24/36 \\
1.1 & 0.2   & 0 & 0 & \warncell{9} & 0/36 \\
1.1 & 1     & 0 & 0 & \warncell{9} & 0/36 \\
\bottomrule
\end{tabular}
\end{table}

With 10\% resistance mismatch and $r=0.05$, six candidates satisfy the 1 A
limit for every state in the box, while three remain unresolved
(Table~\ref{tab:certificate-informativeness}). The selected retained candidate
has $c=s=0$; its peak-current envelope is $[0.8285,0.9667]$ A
(rounded outward). Its regret against the best candidate in the same
retained set is at most $0.417$ in the dimensionless score. The tighter
limits, 0.25 and 0.5 A, certify all nine candidates infeasible for
$r\le0.2$; the 2 A limit retains all nine for $r\le0.2$. At $r=1$, every
limit is unresolved. Thus validity alone does not imply an informative
decision, and lack of certification does not imply infeasibility.

An independent adaptive ODE integration agrees with the nominal matrix flow
within $1.6\times10^{-12}$ in the reported physical coordinates. All 11,529
candidate--vertex diagnostic rollouts satisfy the asserted envelopes,
orderings, feasibility implications, and regret bounds within $10^{-9}$
floating-point checking tolerance. The interval construction, not these
finite tests, establishes the box-wide result. The certificate covers the
sampled horizon only, not intersample overloads or repeated execution.
\FloatBarrier
